\begin{table}[H]
    \centering
    \caption{Interval \textit{Bootstrap} 95\% untuk Estimasi Utama Thorndike Case III (Delapan Subtes)}
    \label{tab:bootstrap}
    \begin{spacing}{1.0}
    \begin{tabular}{lcccc}
        \toprule
        Statistik & Estimasi & Rerata \textit{Bootstrap} & $SE$ \textit{Bootstrap} & Interval Persentil 95\% \\
        \midrule
        $S_Z$ populasi & 67,72 & 67,73 & 0,55 & [66,63; 68,79] \\
        $\lambda^2$ & 0,582 & 0,586 & 0,058 & [0,479; 0,700] \\
        $\bar{R}$ terkoreksi & 0,425 & 0,425 & 0,009 & [0,407; 0,442] \\
        \bottomrule
    \end{tabular}
    \end{spacing}
    \vspace{2mm}
    \begin{flushleft}
    \footnotesize\textit{Catatan.} Interval dihitung dengan \textit{bootstrap} nonparametrik (penarikan ulang entri sampel teramati) sebanyak 2.000 replikasi, 2.000 di antaranya konvergen. Interval hanya menangkap ketidakpastian pengambilan sampel dan mengandaikan $SD$ nasional serta spesifikasi Case III benar (lihat Tabel~\ref{tab:sens_sd}).
    \end{flushleft}
\end{table}
